\newpage
\section{Conclusions}

\subsection{Keyfinding?}

%Considered together, the three BET parameters show that pyrolysis temperature and subsequent treatment strongly influenced the pore structure of the French biochars. Increasing pyrolysis temperature alone resulted in higher specific surface area and pore volume, while the mean pore diameter decreased. This indicates that progressive devolatilization during pyrolysis exposed and developed a finer internal pore network rather than merely generating larger voids.

%The subsequent treatment produced an even stronger modification of the pore structure. Among the investigated precursors, the biochar initially pyrolyzed at 500,$^\circ$C showed the most favorable final BET characteristics, combining the highest specific surface area of 146.09~m$^2$,g$^{-1}$ and the highest pore volume of 0.095~cm$^3$,g$^{-1}$ with the smallest mean pore diameter of 2.603~nm. This suggests that 500,$^\circ$C provided a favorable balance between sufficient initial carbonization and remaining reactivity toward steam activation. In comparison, the 400,$^\circ$C precursor was less carbonized and likely underwent a larger contribution from additional devolatilization during activation, whereas the 600,$^\circ$C precursor was already more strongly carbonized and therefore less susceptible to further pore development.

%Nevertheless, the final BET values represent the combined effects of steam activation and subsequent sulfur loading. Since BET analysis was performed after sulfur loading, the increase in surface area and pore volume cannot be attributed quantitatively to steam activation alone. Sulfur deposited within the pore system may also have occupied part of the available pore volume or partially blocked pore entrances, thereby influencing the surface subsequently accessible to nitrogen during BET measurement. A separate BET analysis directly after activation and before sulfur loading would therefore be required to distinguish the effect of activation from the effect of sulfur deposition.


%Nevertheless, the values cannot be compared solely on the basis of the starting material, as the activation conditions also differed between the studies. Furthermore, in this study, the BET measurements were performed both after steam activation and after sulfur loading, so that the measured surface area reflects the final, sulfur-loaded material and not the biochar immediately after activation.


\subsection{Outlook}

- %bet anlysis cuman buat biochar yang aactivated doang
- %SEM dilaakuin
- %sulfur loading jangan pake kapas lagi, mungkin bisa caari alternativ lain
- %buat activationn pake steam mungkin bisa bikin leitung buat supply waater dengan tekanan yang pas tau sesuai
- %Terus mungkin bisa minin temperatur activationnya dingka 850 (paling bagus sesuai literatur)
- %A direct comparison on one digestate, with BET run on both activated chars before sulfur loading, would be needed to turn this reading into a quantitative result.
- % untuk plausibilas pengaplikasiianny, mungkin bisa dilakukan secara langsung, seperti nanam bunga untuk soil amandment ataau pun sulufr loading (H2S) removl tapi lebih proper














The results of this study


Several aspects can be investigated further in order to better understand and optimize the treatment process.

The results of this study demonstrate that activation and sulfur loading substantially enhance the specific surface area of French biochar samples compared to their pristine counterparts, with the most pronounced effect observed at 500°C. However, several aspects remain open for further investigation and could guide future research in this field.


%The results of this work demonstrate that slaughterhouse-digestate-derived biochar can be modified through steam activation and subsequent sulfur loading. However, 
%A particularly important step for future work would be to characterize the biochar at each individual treatment stage. In the present study, BET analysis was performed on the pristine biochar and on the final activated, sulfur-loaded material. Therefore, the effects of steam activation and sulfur loading on the pore structure could not be separated quantitatively. Future experiments should include BET measurements directly after steam activation and before sulfur loading. In addition, both investigated water doses should be characterized, since the lower dose of 0.1~mL\,g$^{-1}$ resulted in higher sulfur-loading capacities, while BET measurements were only performed for the standard activation condition of 0.5~mL\,g$^{-1}$. This would allow the relationship between activation intensity, pore development and sulfur-loading capacity to be evaluated more directly.

%Further optimization of the activation process could include systematic variation of the activation temperature, duration and steam supply. The present results indicate that the biochar initially pyrolyzed at 500~$^\circ$C was particularly promising in terms of pore development, reaching the highest measured specific surface area of 146.09~m$^2$/g after activation and sulfur loading. Future experiments could therefore focus on this precursor while varying the activation conditions in order to determine whether a higher accessible surface area and pore volume can be obtained without excessive carbon loss.

%The sulfur-loading experiments should also be extended using controlled H$_2$S concentrations and gas flow rates. Breakthrough experiments would allow the actual H$_2$S adsorption capacity and the time-dependent performance of the biochar to be determined more accurately than gravimetric mass gain alone. Additional surface characterization could further clarify the chemical form and distribution of sulfur retained on the biochar.

%Finally, the applicability of the produced biochar should be investigated under conditions closer to practical use. Considering its ability to retain sulfur-containing species, application in H$_2$S removal, for example during biogas purification, represents a particularly relevant direction for further research. Comparison with conventional activated carbons, together with investigations of regeneration, long-term stability and economic feasibility, would be necessary to evaluate whether slaughterhouse digestate can serve as a technically and environmentally viable precursor for adsorbent production.













